\section{Conclusion}
\label{sec-conclude}

An ontology's hierarchy is a better filter for vector search than an attribute the
index cannot see.
Postorder intervals in the store's B+ tree make a class's members a few range scans and
membership a binary search; a vector search can then test the filter per row, or
resolve it and score the members exactly, and a cost model with one threshold picks the
cheaper.
On four ontologies this answers $\mathrm{near}(q,c,k)$ with recall 0.999 and a median of
26\,$\mu$s, where FAISS indexes with ID selectors return recall 0.61 and 0.64 and lose
almost every answer for a query outside a small class, and where an exact FAISS scan is
33 times slower for classes of at most ten members.
No sorted order of the terms can do the same: every order separates some neighbors by
$(N{-}1)/D$ positions, and the best we found holds at most 56\% of a term's neighbors in
a window of 200.

\stitle{Limitations}.
Our IVF and HNSW are 2 to 4 times slower than FAISS's when called one query at a time and
take longer to build, the price of keeping lists in the triple tree; the vectors are held
in memory and the indexes are rebuilt, not updated, when they change; the labels are
rebuilt when the hierarchy changes; the embeddings come from one small model that
embeds names only; masked HNSW is slow for classes of a few hundred members and under,
which the rule avoids and a denser graph, as ACORN builds~\cite{patel2024}, might fix;
and for classes of 1,001 to 5,000 members an exact FAISS scan with a precomputed member
bitmap is faster than our rule.
